% titel: Blatt 8, A4: Approximation eines Integrals
% Teil II: der durchgerechnete Weg.

\subsection*{Gegeben und gesucht}
a) Erklären, was der Python-Code tut (Gauss-Dichte-Integral). b) Den Code so
ändern, dass die Approximation besser wird.

\subsection*{a) Was der Code macht}
\emph{Gedankenschritt:} $f$ ist die Gauss-Dichte von $\mathcal N(0,1)$. Der
Code summiert Rechtecke der Breite $h=\frac1{100}$ über den Bereich
$[-1{,}96,1{,}96]$ (Stützstellen links):
\begin{equation*}
  S=\sum_{k=-196}^{196} f\!\left(\frac{k}{100}\right)\frac{1}{100}
  \approx \int_{-1{,}96}^{1{,}96}
  \frac{1}{\sqrt{2\pi}}\mathrm e^{-x^2/2}\,dx .
\end{equation*}
Der Wert ist $\approx0{,}95$, weil $1{,}96$ das $97{,}5\%$-Quantil ist:
\begin{equation*}
  \Phi(1{,}96)-\Phi(-1{,}96)
  =\Phi(1{,}96)-\bigl(1-\Phi(1{,}96)\bigr)
  =2\Phi(1{,}96)-1\approx2\cdot0{,}975-1=0{,}95 .
\end{equation*}

\subsection*{b) Bessere Approximation}
\emph{Gedankenschritt:} Zwei Stellschrauben: (1) feineres $h$ verkleinert
den Fehler O(h) der Rechteckregel, (2) eine bessere Quadratur nutzt die
Information pro Stück besser.

\emph{1. Feineres $h$:} derselbe Bereich mit $h=10^{-3}$:
\begin{verbatim}
import math
def f(x): return (1/math.sqrt(2*math.pi))*math.exp(-x**2/2)
S = sum(f(k/1000)*(1/1000) for k in range(-1960, 1961))
print(S)
\end{verbatim}

\emph{2. bessere Regel (Mittelpunkt)} mit $h=10^{-3}$:
\begin{verbatim}
# Mittelpunktregel: Stützstelle in der Mitte jedes Teilintervalls
S = sum(f((k+0.5)/1000)*(1/1000) for k in range(-1960, 1960))
print(S)
\end{verbatim}

\begin{ergebnis}
Der Code approximiert
\[
  \int_{-1{,}96}^{1{,}96}\frac1{\sqrt{2\pi}}\mathrm e^{-x^2/2}\,dx
  \approx0{,}95000 .
\]
Verbesserung: Schrittweite verkleinern (z.B.\ $h=10^{-3}$) und eine höhere
Quadraturregel (Mittelpunkt/Trapez/Simpson) wählen.
\end{ergebnis}

\begin{kontrolle}
\textbf{Numerisch (verifiziert):} Die Links-Rechteck-Summe mit $h=10^{-2}$
liefert $\approx0{,}95059$, die Mittelpunkt-Regel mit gleichem $h$
$\approx0{,}9500052$, das wahre Integral $\approx0{,}9500042$. Die
Mittelpunkt-Regel ist also deutlich genauer als die (einseitig
ausgerichtete) Rechteckregel -- genau die $O(h)$- gegen
$O(h^2)$-Fehlerordnung. \checkmark

\textbf{Plausibilität:} für den Bereich einer Faktor-$\sigma$ ist der
Integralwert $0{,}95$ (90\%- vs.\ 95\%-Konfidenzbereich von
$\pm1{,}645$ zu $\pm1{,}96$). \checkmark

\textbf{Grenzfall:} für $b\to\infty$ (ganze reelle Achse) ginge das
Integral gegen $1$; hier ist der Bereich endlich $[-1{,}96,1{,}96]$,
deshalb $0{,}95$. \checkmark
\end{kontrolle}

\begin{merksatz}
Eine Riemann-Summe $\sum f(x_i)\,h$ approximiert ein Integral mit Fehler
$O(h)$. Genauer wird es durch kleineres $h$ oder durch bessere
Quadratur-Regeln (Mittelpunkt/Trapez $O(h^2)$, Simpson $O(h^4)$). Der
Bereich $[-1{,}96,1{,}96]$ ist hier fest; für $+\infty$ müsste man ihn
erweitern.
\end{merksatz}
